\chapter{Trabalhos Relacionados}
\label{sec:relacionados}

Diversos trabalhos na literatura abordaram a utilização de veículos aéreos não tripulados (VANTs) como ferramentas para coleta de dados em redes de sensores da Internet das Coisas (IoT), com foco em otimização de trajetória, eficiência energética e minimização da idade da informação (AoI).

Wei et al.~\cite{wei2022survey} apresentam um extenso levantamento sobre coleta de dados assistida por VANTs em IoT, abordando os principais desafios e tecnologias associadas, como clustering de sensores, modos de coleta (voo pairado, sobrevoo e híbrido) e planejamento de trajetória conjunto com alocação de recursos. Os autores destacaram que a flexibilidade dos VANTs permite melhorar significativamente a cobertura e eficiência da coleta de dados em ambientes complexos. Além disso, o trabalho propõe uma taxonomia das abordagens existentes e discute algoritmos baseados em teoria dos grafos, teoria da otimização e inteligência artificial para resolver o problema.

Messaoudi et al.~\cite{messaoudi2023survey} também forneceram uma revisão abrangente dos esquemas de coleta de dados baseados em VANTs, com foco em desafios como consumo energético, atraso na coleta, limitação de mobilidade e segurança. O artigo propôs uma categorização das aplicações e técnicas, discutindo ainda abordagens energéticas como redes com carregamento sem fio e estratégias para minimizar a AoI. Os autores enfatizaram a importância de integrar VANTs de forma inteligente para garantir eficiência e longevidade da rede IoT, além de sugerirem direções futuras promissoras, como o uso de VANTs energizados por comunicação e integração com redes 5G/6G.

Liu et al.~\cite{liu2018uavWSN} investigaram o desempenho da coleta de dados em redes de sensores sem fio (WSNs) com apoio de VANTs, especialmente em cenários sem infraestrutura de comunicação. O artigo propôs um modelo onde a região de cobertura é dividida em células e diferentes estratégias de trajeto são projetadas para um ou múltiplos VANTs. Os autores derivaram a capacidade por nó (\textit{per-node capacity}), que depende do número de células, da altura dos VANTs, da densidade de sensores e da capacidade energética das aeronaves. Concluiu-se que o uso de múltiplos VANTs aumenta significativamente essa capacidade em comparação com apenas um VANT. Além disso, o artigo determina o número ótimo de células para maximizar o desempenho da rede e valida os resultados teóricos com simulações. O trabalho contribui para o entendimento do impacto de diferentes parâmetros na eficiência da coleta de dados, oferecendo orientações práticas para o planejamento de missões VANT em WSNs.

Zhang et al.~\cite{zhang2018drlrvc} propuseram uma estrutura baseada em aprendizado profundo para otimizar a coleta de dados por veículos autônomos em cidades inteligentes, com foco especial em eficiência energética. O trabalho introduziu o framework \textit{DRL-RVC}, que utiliza redes neurais convolucionais para extrair características relevantes (como distribuição de amostras e fluxo de tráfego), e redes Q profundas para controlar VANTs e veículos autônomos de recarga. Os VANTs realizam a coleta de dados priorizando áreas de maior relevância, enquanto veículos autônomos terrestres são mobilizados para oferecer recarga, contornando restrições urbanas. Os resultados de simulação mostraram alta efetividade e robustez da abordagem proposta.

Yang et al.~\cite{yang2020surveyUAVIoT} realizam uma análise abrangente das principais questões na coleta de dados assistida por VANTs em ambientes de IoT, abordando desde o impacto do posicionamento dos nós sensores até o planejamento de trajetória e navegação autônoma. O estudo destacou que uma implantação eficiente dos sensores, aliada à otimização de parâmetros de voo (como velocidade e altitude) e a técnicas de compressão, clustering e classificação de dados, é fundamental para maximizar a eficiência energética e a quantidade de dados coletados. Os autores também enfatizaram o papel de algoritmos de planejamento de rota e desvio inteligente de obstáculos, sugerindo que abordagens inspiradas em aprendizado por reforço podem melhorar significativamente o desempenho de VANTs na coleta em larga escala, especialmente em terrenos complexos e de difícil acesso.

A Tabela~\ref{tab:comparativo} sintetiza a comparação entre os trabalhos revisados e a proposta deste trabalho. As duas últimas colunas indicam se cada trabalho trata a energia e a AoI como \emph{objetivos de otimização}, ou seja, grandezas efetivamente minimizadas ou maximizadas por um modelo, e não apenas como temas discutidos de forma qualitativa. Sob esse critério, os \textit{surveys} de Wei et al., Messaoudi et al. e Yang et al. mapeiam desafios de consumo energético e de atualidade dos dados, mas não formulam um modelo que os otimize. O trabalho analítico de Liu et al. concentra-se em maximizar a capacidade por nó, sem tratar energia nem AoI como objetivo. Zhang et al. otimizam a eficiência energética por aprendizado profundo, porém não consideram a AoI. Constata-se, assim, que nenhum trabalho anterior otimiza AoI e energia conjuntamente.

Este trabalho preenche essa lacuna ao formular o problema de coleta de dados como um programa linear inteiro misto (MILP) com dois objetivos conflitantes: maximizar a AoI coletada pelos sensores e minimizar o consumo energético do VANT. A abordagem lexicográfica empregada garante soluções ótimas exatas, ao passo que o uso de parâmetros físicos reais de dois modelos de drone (DJI Matrice 300 RTK e Parrot Anafi USA) confere aderência prática ao modelo. Diferentemente dos trabalhos que consideram múltiplos VANTs, optou-se pelo cenário com VANT único, o que permite isolar o impacto das decisões de trajetória e energia sobre a qualidade da coleta de forma controlada.

\begin{table}[htbp]
  \centering
  \caption{Comparativo entre trabalhos relacionados e a proposta deste trabalho. As colunas ``Otimiza energia'' e ``Otimiza AoI'' indicam se o trabalho adota cada grandeza como objetivo de otimização, e não apenas como tema de discussão.}
  \label{tab:comparativo}
  \small
  \adjustbox{max width=\linewidth}{%
  \begin{tabular}{>{\raggedright\arraybackslash}p{4.3cm}
                  >{\centering\arraybackslash}p{1.1cm}
                  >{\centering\arraybackslash}p{1.7cm}
                  >{\raggedright\arraybackslash}p{2.6cm}
                  >{\centering\arraybackslash}p{1.7cm}
                  >{\centering\arraybackslash}p{1.4cm}}
    \toprule
    \textbf{Trabalho} & \textbf{Ano} & \textbf{N° VANTs} & \textbf{Método} & \textbf{Otimiza energia} & \textbf{Otimiza AoI} \\
    \midrule
    Wei et al.~\cite{wei2022survey} & 2022 & Múltiplos & Revisão & Não & Não \\
    Messaoudi et al.~\cite{messaoudi2023survey} & 2023 & Múltiplos & Revisão & Não & Não \\
    Liu et al.~\cite{liu2018uavWSN} & 2018 & Múltiplos & Analítico & Não & Não \\
    Zhang et al.~\cite{zhang2018drlrvc} & 2018 & Múltiplos & Aprendizado profundo & Sim & Não \\
    Yang et al.~\cite{yang2020surveyUAVIoT} & 2020 & Múltiplos & Revisão & Não & Não \\
    \midrule
    \textbf{Este trabalho} & \textbf{2026} & \textbf{Único} & \textbf{MILP (exato)} & \textbf{Sim} & \textbf{Sim} \\
    \bottomrule
  \end{tabular}}
\end{table}
